\section{Proof of Theorem~\ref{thm:main}}\label{sec:proof}

\subsection{The criterion}
We prove Proposition~\ref{prop:criterion}. Its proof shows where each of the three properties of
$Q_N$ enters: the nonvanishing and the smallness are used once at the end, and the divisibility by prime
squares is what makes the choice of the prime free.

\begin{proof}[Proof of Proposition~\ref{prop:criterion}]
Let $a/b$ be a rational number with $b$ large, put $L=\log b$, $U=\kappa L$ and $V=(\kappa-\omega)L$, and
suppose $|\xi-a/b|\le b^{-\rho\kappa}$. By the prime number theorem, $(1-\delta)t\le\theta(t)\le(1+\delta)t$
for all large $t$, and $V\to\infty$ with $b$. If every prime $p\in(V,U]$ had $p^2\mid b$, then
$2(\theta(U)-\theta(V))\le\log b=L$; but
\[
 2\bigl(\theta(U)-\theta(V)\bigr)\ge2L\bigl((1-\delta)\kappa-(1+\delta)(\kappa-\omega)\bigr)
 =2L\bigl(\omega-\delta(2\kappa-\omega)\bigr)>L ,
\]
by the fourth hypothesis. So there are primes in $(V,U]$ whose squares do not divide $b$. Let $p$ be the
\emph{largest} one. Every prime $q\in(p,U]$ has $q^2\mid b$, and $q\le U<2V<2p$ because $2\omega<\kappa$.
Let $B$ be the product of these $q^2$; then $B\mid b$ and $\log B=2(\theta(U)-\theta(p))$.

Every term $c_ia^ib^{dp-i}$ of the integer $b^{dp}Q_p(a/b)=\sum_ic_ia^ib^{dp-i}$ is divisible by $q^{2dp}$ for each
of these $q$: $q^{2i}$ divides $c_i$ by hypothesis, and $q^2\mid b$ gives $q^{2(dp-i)}\mid b^{dp-i}$. So $B^{dp}$
divides the integer. Since $Q_p(a/b)\ne0$ (as $p^2\nmid b$), $b^{dp}Q_p(a/b)/B^{dp}$ is a nonzero integer, and
\begin{equation}\label{eq:recycled}
 1\le\Bigl(\frac bB\Bigr)^{dp}\,|Q_p(a/b)| .
\end{equation}
Since $\rho p\le\rho U=\rho\kappa L$, we have $|a/b-\xi|\le e^{-\rho p}$, so $|Q_p(a/b)|\le e^{-\eta p^2}$ and
\eqref{eq:recycled} gives $\eta p^2\le dp\bigl(L-2(\theta(U)-\theta(p))\bigr)$, that is,
\[
 \eta p\le d\bigl(L-2(\theta(U)-\theta(p))\bigr)\le d\bigl(L-2U+2p+4\delta U\bigr),
\]
using $\theta(U)\ge(1-\delta)U$ and $\theta(p)\le(1+\delta)p\le(1+\delta)U$ in the form
$\theta(U)-\theta(p)\ge U-p-2\delta U$. Hence $(\eta-2d)p\le dL-2dU+4d\delta U$. As $\eta<2d$ and $p\le U$,
the left side is at least $(\eta-2d)U$, so $\eta U\le dL+4d\delta U$, that is, $\eta\kappa\le d+4d\delta\kappa$,
which contradicts the last hypothesis.
\end{proof}

The choice of the largest admissible prime is what makes the selection free: the primes of the window
that are not admissible have their squares in $b$, and these squares are recovered from the coefficients
through the divisibility hypothesis. A cruder choice, of any prime in a short interval that does not
divide $b$, loses a factor. The free lower end $V=(\kappa-\omega)L$ is what allows $\kappa<2$: the condition $q\le U<2p$ for the recycled
primes needs only $U<2V$.

\subsection{The hypotheses for the design $D$}
Take the design $D=(58,9,3,5,3)$ and $d=53$ (the degree of $Q_N$ is at most $53N$, and the polynomials
$Q_p$ of the proposition are the $Q_N$ with $N=p$), and
\begin{gather*}
 \rho=\Rho,\quad \eta=\EtaTot,\quad \omega=\tfrac34,\quad \delta=2.9\cdot10^{-7},\\
 \kappa=\frac{\Bbank}{\rho}=\Kappa .
\end{gather*}
We check the hypotheses.
\begin{itemize}
\item $\deg Q_N\le53N$ by \eqref{eq:QN} and Lemma~\ref{lem:degree}.
\item The nonvanishing: by Proposition~\ref{prop:nonvanishing}, $\Delta_p(a/b)\ne0$ for large $p$ with
$p^2\nmid b$, hence $Q_p(a/b)\ne0$.
\item The divisibility: by Lemma~\ref{lem:recycling}, $q^{2i}$ divides the $i$-th coefficient of $Q_p$ for all
primes $p<q\le2p$.
\item The smallness: by Theorem~\ref{thm:real}, Proposition~\ref{prop:denominators} and $\log(53N)!=o(N^2)$,
for $|x-G|\le e^{-\rho N}$ and large $N$,
\[
 |Q_N(x)|\le\exp\bigl(-(\TDelta\ldots-6395.7087\ldots-3\varepsilon)N^2\bigr)\le e^{-\eta N^2}
\]
since $\eta=\EtaTot<\EtaExact\ldots$ is rounded down.
\item The numerical conditions: $\rho\kappa=\Bbank$; $0<\eta<2d=106$; $\omega=3/4\le1$ and $2\omega=1.5<\kappa=\Kappa$;
$2(\omega-\delta(2\kappa-\omega))=1.49999\ldots>1$; and
$d+4d\delta\kappa=\PayL\ldots<\eta\kappa=\PayR\ldots$.
\end{itemize}
Proposition~\ref{prop:criterion} therefore gives $|G-a/b|>b^{-\rho\kappa}=b^{-\Bbank}$ for all integers
$a$ and all sufficiently large $b$. This proves Theorem~\ref{thm:main}. The last condition is
equivalent to $\Bbank>x=d\rho/(\eta-4d\delta)=\Bx$, and in the limit $\delta\to0$ the argument gives
the exponent $d\rho/\eta=\Blimit\ldots$ for the constants above; the integer $\Bbank$ is the one for
which all conditions are verified with rational constants.

\section{Remarks on the method}\label{sec:remarks}

The exponent $d\rho/\eta$ has three ingredients. The degree $d=n-g$ and the perturbation exponent
$\rho=n_0c_2+w$ come from the construction and from the quadrature: of $\rho=\Rho$, the cost of fixed
columns $n_0c_2$ is $\NCtwo$ and the logarithmic size of the weights is $52.79$. The decay rate $\eta$ is
a small difference of two large numbers, the energy at the real place, $\TDelta$, and the exponent of the
denominators, $6395.71$, and the rank deficiencies of Theorem~\ref{thm:content} contribute $19.14$ of the
$\EtaTot$. Since the exponent is inversely proportional to $\eta$, an additional gain of $0.01$ in the area
of $\kappa$ lowers it by about $0.22$.

All four theorems of Section~\ref{sec:content} are lower bounds for the nullity, and the proofs give the
kernel vectors explicitly. They are proved for every design, and so are the other arguments of the paper; a
different design changes $d$, $\rho$ and $\eta$ together, and we make no claim that $D$ is optimal. In
particular the bound $\kappa$ of Theorem~\ref{thm:content} does not use any relation between different
families, and at a prime where windows overlap the exact nullity is larger than each bound. An upper bound
for the nullity, or a proof that the kernels of two families are independent, would increase $\eta$ without
changing anything else.

\section{Formalization}\label{sec:formal}

All results of this paper are formalized in Lean~4 \cite{Lean4} with Mathlib \cite{Mathlib}, on top of
OpenAI's formalization of the irrationality of $G$ \cite{OpenAI2026}. The development is available at
\begin{center}
\FormalizationLink.
\end{center}
The statement of Theorem~\ref{thm:main} in \lean{Challenge.lean} uses only Mathlib, with Catalan's constant
written as the series $\sum_j(-1)^j/(2j+1)^2$; it is the theorem
\begin{center}
\texttt{CatalanIrrationalityMeasure.catalan\_irrationalityMeasure\_\Bbank}
\end{center}
and its proof is \lean{Solution.lean}. The solution has no \lean{sorry} and no \lean{native\_decide}, and the
command \lean{\#print axioms} reports only the three standard axioms \lean{propext}, \lean{Classical.choice} and
\lean{Quot.sound}. The file \lean{comparator.json} is a configuration for \lean{lake comparator}
\cite{Comparator} stating that \lean{Solution} proves the statement of \lean{Challenge} with these axioms
only. The prime number theorem is taken from the PNT+ project \cite{PNTPlus}. The parts of OpenAI's
formalization that are used were converted to Lean's module system and ported to a recent Lean release; the
repository's \texttt{README} and \texttt{PORTING.md} describe the changes.

\subsection{Module layout}
Everything specific to this paper is in the directory \lean{GenD} of \lean{CatalanMeasure}. Its design is that of the paper:
a design is data, the generic modules prove every statement for an arbitrary design, and the instance
files contain numerals only.
\begin{description}[leftmargin=1.5em,style=nextline]
\item[Generic layer] \lean{PL} (piecewise linear profiles, their areas) and \lean{PrimeGrowth} (the
weighted prime number theorem \eqref{eq:weighted-pnt} for ramp profiles; one exponent per prime);
\lean{Budget} (the scalar inequalities of Lemma~\ref{lem:budget2} and its analogue for the first case, and the
conditions of the criterion), \lean{PayW} and \lean{Recycled} (Proposition~\ref{prop:criterion} with the free
window end); \lean{Final} (the structure \lean{DesignData} and the theorem \lean{measure\_of\_design}, which
contain no numeral).
\item[Denominators] \lean{OddProfile} (the valuations $\ell_N(p)$ for the lower and upper pools),
\lean{TwoAdic} (the bound \eqref{eq:A2}), \lean{DenomExp}, \lean{DenomGrowth}, \lean{DenomGrowth2},
\lean{Denominator} and \lean{Normalizer} (the common denominator, $Q_N$, the divisibility of
Lemma~\ref{lem:recycling}).
\item[Rank deficiencies] \lean{Interface} (the nullity, the four families and \lean{famKap\_le\_nullity}),
\lean{RankCommon}, \lean{RankK6}, \lean{RankDigits}, \lean{RankFam}, \lean{RankFilt}, \lean{RankSeq},
\lean{RankBlock} (the shared layers: Chebyshev reflection, identity \eqref{eq:K6}, the digit formula in
filtered form, the block framework), and \lean{RankT}, \lean{RankU}, \lean{RankUCore}, \lean{RankU6},
\lean{RankS} for Theorems~\ref{thm:T}--\ref{thm:S}; \lean{KapHelpers} turns the comparison with a window list
into linear arithmetic.
\item[Nonvanishing] \lean{Nonvanishing} (Proposition~\ref{prop:nonvanishing} from the three base matrices),
and the residue-table certificates \lean{CertMod}, \lean{ResBase}, \lean{ResCheck}, \lean{ResInv},
\lean{ResRec}.
\item[Real place] \lean{Analytic/} (28 modules): the constants and fields of a design (\lean{Const},
\lean{Fields}, \lean{CommonField}), the first-case and second-case contracts (\lean{First},
\lean{FirstGlue}, \lean{Second}, \lean{Glue}), the grouped products over the fixed points (\lean{GroupedProduct},
\lean{MarkedSheet}, \lean{MarkedSheetOn}, \lean{MarkedMixedMajorant}, \lean{MarkedSecond}), the prefactor,
the perturbation by the quadrature nodes (\lean{Prefactor}, \lean{AtomBase}, \lean{WeightedAtom},
\lean{DeltaSmall}, \lean{DeltaChain}), and the nodes and the weight bounds for any number of nodes
(\lean{GapNodes}, \lean{GapNodal}, \lean{GapWeights}, \lean{FullWeights}, \lean{GapCircleGen},
\lean{CircleCover}). \lean{Cert/} (6 modules) is the generic cell machinery: the evaluator of logarithms
(\lean{WLog}, \lean{WLogSpec}), the residual field with a weight (\lean{ResW}), the shifted tables
(\lean{PolyShift}) and the assembly of cells and blocks (\lean{CellLoc}, \lean{Assembly}).
\item[The design $D$] \lean{D58/} holds the instance: \lean{Budget} (the constants of the certificate),
\lean{Ramps}, \lean{Kap}, \lean{OddProfile}, \lean{TwoAdic}, \lean{Normalizer}, \lean{Pay}, \lean{Final}
(the design, the window list, the profile $f$, $\mathsf A_2=2673/8$, the growth of the denominator and
the conditions of the criterion), \lean{Circle} (the 42-piece cover), \lean{FirstTrial},
\lean{FirstTrialCert}, \lean{FirstBridge} and \lean{FirstCert/} (the first case), \lean{Cert/} (the second
case: trial data, the covers of $X$, $Y$ and $R$ in blocks, the norms, the near-field cap),
\lean{SecondTrial}, \lean{C3/} (the three inverses modulo $1009$), \lean{Certified} and \lean{Main}, which
proves \lean{catalan\_irrationalityMeasure\_\Bbank} in the namespace \lean{CatalanMeasure.GenD.D58}.
\end{description}
The public theorem in \lean{Solution.lean} is a one-line wrapper of that theorem. The modules of
OpenAI's formalization and the generic valuation theorems of the construction
(\lean{CatalanMeasure/Parametric}) are used as they are.

\subsection{Certificates}
The proof combines conventional formal proofs with numerical certificates checked by the Lean kernel:
the inverses of the base matrices modulo $1009$, the cover of Lemma~\ref{lem:quadrature}, the trial
sequences of Section~\ref{ssec:trials} and the covers of Section~\ref{ssec:cells}. All are exact rational or
integer computations, decided by the kernel without compiled code. The interval arithmetic used to
generate the trial data and the covers \cite{Arb} is not trusted: Lean rechecks every inequality. Four
devices keep the certificates small enough for the kernel.
\begin{itemize}
\item \emph{Shared Taylor bounds.} The cells of a cover are grouped in blocks. For each block the
derivatives of the polynomial part of the field are evaluated once, exactly, at the centre of the block,
together with one bound for the remainder of the expansion; every cell of the block obtains its expansion
by an exact shift, instead of repeating the evaluation of a polynomial with $\Kcoef$ coefficients.
\item \emph{Accuracy per term from a shared table of nodes.} The logarithms are evaluated by the series
$\log\frac{1+x}{1-x}=2x\sum_i x^{2i}/(2i+1)$ with an explicit geometric tail, and the number of terms is
chosen for each logarithmic term separately from a table of nodes shared by all the cells, instead of
one accuracy for the whole sum.
\item \emph{Residue recurrences.} The tables of the moments $M(i,j)$ and $Z(i,j)$ modulo $1009$ are not
computed from the rational numbers. They are generated by the one-step recurrences that define these
numbers, once proved correct, and the kernel checks the recurrence, the filtered rows against the
row polynomials and the products of the matrices with their claimed inverses, row by row, in natural
numbers.
\item \emph{$2$-adic normalisation.} Where a recurrence produces integers divisible by large powers of
$2$, as the Chebyshev and Clenshaw recurrences do, the tables store the quotients by these powers, with
one lemma relating the two; this also avoids integers that the kernel's hashing of literals treats as
equal.
\end{itemize}

\section*{Disclosure of AI use}

The mathematics, the Lean formalization, the numerical computations and certificate generators, the
literature search and the text of this paper were produced by AI systems at the request of the author, who
posed the problem, directed the work, and is responsible for this paper. By component: Claude
(Opus~5.5 and Sonnet~5.5, Anthropic) produced the treatment of an arbitrary design (the degree, the
denominators at all primes and the nonvanishing, Section~\ref{sec:setup}), the digit formula and the rank
deficiencies of Theorems~\ref{thm:T}, \ref{thm:U}, \ref{thm:U6} and~\ref{thm:S}, the passage from the
nullity to the denominators, the criterion with a free lower end of the window, the perturbation estimate
by quadrature with an odd number of nodes, the trial sequences with complex terms and the verification of
Section~\ref{ssec:cells}, the choice of the design and the final assembly, together with the corresponding
Lean proofs and the generic layer of Section~\ref{sec:formal}. OpenAI Codex produced the energy estimates with
fixed columns and the fields of the first case, the identity of Andr\'eief type with fixed points, the
bound for the weights of Lemma~\ref{lem:quadrature}, the divisibility of the coefficients by prime squares
and the choice of the largest admissible prime, with the corresponding Lean proofs for one particular
choice of the parameters, which Claude generalised. The underlying irrationality proof and its
formalization are OpenAI's \cite{OpenAI2026}.

\subsection*{Acknowledgements}
We thank the authors of Mathlib and of the PNT+ project, and OpenAI for making its formalization available.
